% Данный файл распространяется под лицензией CC BY 4.0. Текст лицензии размещён на https://creativecommons.org/licenses/by/4.0/
% (c) Кафедра системного программирования, 2025

\documentclass[a4paper]{article}

\usepackage[a4paper, top=8mm, bottom=8mm, left=8mm, right=8mm]{geometry}

\usepackage{polyglossia}
\setdefaultlanguage[babelshorthands=true]{russian}
\setotherlanguage{english}

\usepackage{fontspec}
\setmainfont{FreeSerif}
\newfontfamily{\russianfonttt}[Scale=0.7]{DejaVuSansMono}

\usepackage[tiny, compact]{titlesec}

\usepackage{titling}
\setlength{\droptitle}{-1cm}
\pretitle{\begin{center}\begin{bfseries}\Large}
\posttitle{\par\end{bfseries}\end{center}}
\preauthor{\begin{center}\normalsize}
\postauthor{\par\end{center}\vspace{-1.8cm}}

\usepackage{hyperref}
\usepackage{bookmark}
\usepackage{csquotes}
\usepackage{float}
\usepackage[useregional]{datetime2}
% Сюда можно дописывать нужные вам пакеты.

\title{UCFS: Доработка репозитория}

\author{Щелокаев Григорий Сергеевич}

% Дата много места занимает, её лучше не указывать.
\date{}

\begin{document}

\maketitle

\begin{flushright}
    Группа: \emph{\enquote{25.Б72-мм}}

    Кафедра: \emph{Системного программирования}

    % Степень/звание и должность научника мы лучше вас знаем, их можно не указывать.
    Научный руководитель: \emph{Григорьев Семён Вячеславович}

    Номер семестра практики: \emph{2}
\end{flushright} 


\section{Введение}

Работа выполняется в рамках проекта UCFS\footnote{Ссылка на UCFS: \url{https://github.com/FormalLanguageConstrainedPathQuerying/UCFS} (дата обращения: \DTMdate{2026-09-11}).}. UCFS~--- это инструмент для решения задач анализа графов с дополнительными ограничениями на пути, заданными контекстно-свободной грамматикой какого-либо языка. Данный инструмент реализован преимущественно на языках программирования Kotlin и Java.
На момент начала работы в данном проекте в репозитории не был настроен CI, отсутствовали контроль тестового покрытия и автоматическая публикация jar-файлов.

Данные практики, которые были описаны выше и которые отсутствовали в проекте, являются одними из ключевых и широко используемых способов повышения качества программного продукта. CI/CD обеспечивает автоматизацию сборки, тестирования программы, что приводит к увеличению вероятности обнаружить ошибки в исходном коде. Контроль тестового покрытия позволяет оценить готовность применения продукта в использовании его по назначению. Добавление всех вышеописанных практик в репозиторий в интересах конечных пользователей программного продукта, а также самих разработчиков проекта.

\section{Постановка задачи}

Перед началом работы в проекте совместно с научным руководителем была поставлена общая задача всей работы~--- обеспечение качества проекта. Данная цель, сформулированная в общем виде, декомпозируется на следующие пункты.
\begin{itemize}
    \item Настройка CI в проекте.
    \item Обеспечение контроля тестового покрытия локально и на платформе GitHub.
    \item Настройка публикации jar-файлов с указанием зависимостей программы.
\end{itemize}

В ходе работы в данном репозитории совместно с одним из ведущих разработчиков UCFS было принято решение обновить версию системы сборки проектов Gradle до новейшей на тот момент версии.


\section{Описание предлагаемого решения}


Далее описано решение каждой конкретной задачи, а также требования к её выполнению.
\begin{itemize}
    \item CI. Основные требования: автоматическая сборка проекта, запуск тестов, а также проверка соответствия тестового покрытия установленным порогам при push и pull request на платформе GitHub Actions. Для решения поставленной задачи был использован язык сериализации данных YAML. На нём была написана последовательность действий, которая решает данную проблему. Также был использован action \verb|jacoco-report|~--- \url{https://github.com/Madrapps/jacoco-report} (дата обращения: \DTMdate{2026-09-18}) для составления отчётов о тестовом покрытии в комментариях к соответствующему pull request.
    \item Контроль тестового покрытия. Требования: сбор сведений о тестовом покрытии для следующих метрик: процент протестированных инструкций, веток, строк исходного кода, методов и классов, установление для данных метрик порогов покрытия, а также обеспечение проверки соответствия реальных показателей покрытия требуемым. Для реализации данной задачи был использован плагин JaCoCo~--- \url{https://github.com/jacoco/jacoco} (дата обращения: \DTMdate{2026-09-18}), совместно с научным руководителем были определены пороги покрытия. Сбор сведений о покрытии осложнялся тем, что для модулей проекта solver и generator были как Unit-тесты, так и тесты, расположенные в отдельном модуле test-shared. Решение данной задачи также потребовало создания отдельного файла сборки проекта в корне репозитория. Реализация данного решения была написана на языке Kotlin Script.
    \item CD/публикация пакетов на GitHub. Требования: публикуемый пакет должен быть разбит на jar-файлы модулей с указанием зависимостей проекта в POM-файле. Конфигурация публикации была реализована в вышеупомянутом корневом файле сборки проекта, а также в специальном файле ci-publishing-jar-files.yaml, так, чтобы при выходе новой версии продукта требуемые пакеты выкладывались на GitHub.
\end{itemize}

\section{Ручное тестирование}
На данный момент решение поставленных задач полностью одобрено ведущими разработчиками и влито в форк основного репозитория одного из ключевых разработчиков. Ссылка на закрытый pull request с влитыми в форк изменениями~--- \url{https://github.com/vkutuev/UCFS/pull/4} (дата обращения: \DTMdate{2026-09-11}). Далее представлено ручное тестирование предлагаемого решения посредством запуска Gradle tasks и workflows и сравнения результатов их работы с ожидаемыми в зависимости от состояния проекта.

Проверка выполненных настроек CI/CD, а также контроля тестового покрытия и публикации пакетов проводилась локально и на GitHub Actions в форке~--- \url{https://github.com/GRIGAeo/UCFS} (дата обращения: \DTMdate{2026-09-11}).

Для проверки корректности контроля тестового покрытия из корня проекта были запущены следующие команды.
\begin{itemize}
    \item \verb|./gradlew testCodeCoverageReport|~--- данная команда составляет отчёты о тестовом покрытии для всего проекта в формате .csv, .html и .xml. Если проект по какой-то причине не собирается или не проходят тесты, то команда падает с ошибкой.
    \item \verb|./gradlew jacocoTestCoverageVerification|~--- это команда так же, как и предыдущая, составляет соответствующие отчёты, но помимо этого устанавливает пороги тестового покрытия (т.~е. если реальные проценты по покрытию ниже установленных, сборка закончится с ошибкой). На данный момент проект не соответствует заранее оговорённым с научным руководителем требованиям к порогам тестового покрытия, поэтому данный процесс падает с ошибкой. При занижении порогов тестового покрытия до минимальных значений сборка, запущенная данной командой, корректно завершается.
\end{itemize}

Для проверки корректности настроенного CI/CD использовалась платформа GitHub Actions в данном форке~--- \url{https://github.com/GRIGAeo/UCFS/tree/test-infrastructure} (дата обращения: \DTMdate{2026-09-11}). Далее представлены ссылки на различные сценарии работы workflow \foreignquote{english}{Run tests and test coverage}, в зависимости от того, в каком состоянии находится проект.

\begin{enumerate}
    \item Результат работы workflow, если не проходит какой-то из тестов~--- \url{https://github.com/GRIGAeo/UCFS/actions/runs/30174893883/job/89721972696} (дата обращения: \DTMdate{2026-09-11}). На шаге Run tests в строках 49--54 видно, что тест TestDotParser не прошёл, поэтому весь workflow закончился с ошибкой.
    \item Результат работы workflow, если все тесты прошли, но тестовое покрытие ниже установленных порогов~--- \url{https://github.com/GRIGAeo/UCFS/actions/runs/30175041391/job/89722339260} (дата обращения: \DTMdate{2026-09-11}). На шаге Verify coverage в строках 73--77 видно, что тестовое покрытие проекта ниже установленных порогов, поэтому workflow завершился с ошибкой, а также в PR по данной ссылке \url{https://github.com/GRIGAeo/UCFS/pull/2} (дата обращения: \DTMdate{2026-09-11}) был создан комментарий-таблица с указанием реального покрытия проекта тестами.
    \item Результат работы workflow, если все тесты прошли и тестовое покрытие не ниже установленных порогов~--- \url{https://github.com/GRIGAeo/UCFS/actions/runs/30175355975} (дата обращения: \DTMdate{2026-09-11}). Здесь же видно, что все тесты прошли, покрытие не ниже установленных порогов.
\end{enumerate}

Для проверки корректности конечных jar-файлов со всеми скомпилированными .class-файлами и POM-файла с указанием ссылок на зависимости проекта была локально запущена из корня проекта следующая команда.
\begin{itemize}
    \item \verb|./gradlew publishToMavenLocal|~--- в результате получаются по 2 jar-файла для каждого модуля, содержащих все скомпилированные файлы исходного кода без внешних библиотек, и POM-файл с указанием ссылок на все 6 зависимостей проекта.
\end{itemize}

Для проверки корректности публикации пакетов в очередной раз использовалась платформа GitHub Actions в данном форке \url{https://github.com/GRIGAeo/UCFS} (дата обращения: \DTMdate{2026-09-11}). Далее представлены варианты завершения workflow \foreignquote{english}{Publish package to GitHub Packages}, который запускается только при выходе новых версий проекта.
\begin{enumerate}
    \item Результат работы workflow, если не проходит какой-то из тестов~--- \url{https://github.com/GRIGAeo/UCFS/actions/runs/30455730146/job/90588686710} (дата обращения: \DTMdate{2026-09-11}). На шаге Run tests в строках 70--77 указано, что тест TestDotParser не прошёл, поэтому весь workflow завершился с ошибкой.
    \item Результат работы workflow, если прошли все тесты~--- \url{https://github.com/GRIGAeo/UCFS/actions/runs/30456362241} (дата обращения: \DTMdate{2026-09-11}). Также по данной ссылке \url{https://github.com/GRIGAeo?tab=packages&repo_name=UCFS} (дата обращения: \DTMdate{2026-09-11}) размещены получившиеся скомпилированные файлы двух модулей, удовлетворяющие вышеперечисленным требованиям к публикуемым файлам.
\end{enumerate}

Вывод: все компоненты CI/CD корректно настроены и проверены~--- тесты запускаются, покрытие контролируется, публикация в GitHub Packages работает.

\section{Заключение}

В ходе выполнения работы получены следующие результаты.

\begin{itemize}
    \item После обсуждения с одним из разработчиков инструмента, по совместительству являющимся научным руководителем, были установлены следующие пороги тестового покрытия: INSTRUCTION~--- 0.95, BRANCH~--- 0.8, LINE~--- 0.8, METHOD~--- 0.85, CLASS~--- 0.9.
    \item Результаты тестового покрытия для всего проекта представлены в таблице~\ref{tab:solver-coverage}.
    \item Публикация проекта была корректно настроена в корневом файле сборки проекта, а также в файле настройки CI ci-publishing-jar-files.yaml. Данная конфигурация при новых релизах публикует скомпилированные jar-файлы двух модулей, а также POM-файл с указанием всех внешних зависимостей проекта.
    \item В файлах .github/workflows/ci-publishing-jar-files.yaml, .github/workflows/ci-test-infrastructure.yaml корректно реализованы процессы CI/CD с запуском тестов, проверкой порогов тестового покрытия и публикацией пакетов. 
\end{itemize}

Материалы, иллюстрирующие выполнение работы, отсутствуют.

Ссылка на влитый pull request в репозитории одного из ключевых разработчиков инструмента: \url{https://github.com/vkutuev/UCFS/pull/4} (дата обращения: \DTMdate{2026-09-11}).

Ссылка на fork: \url{https://github.com/GRIGAeo/UCFS} (дата обращения: \DTMdate{2026-09-11}).

Решение одобрено научным руководителем, а также одним из ключевых разработчиков UCFS, предложенные изменения влиты в форк основного репозитория.

\begin{table}[H]
\centering
\begin{tabular}{|l|c|c|}
\hline
Показатель & Покрытие & Порог \\
\hline
INSTRUCTION & 51\% & 95\% \\
BRANCH      & 41\% & 80\% \\
LINE        & 49\% & 80\% \\
METHOD      & 45\% & 85\% \\
CLASS       & 64\% & 90\% \\
\hline
\end{tabular}
\caption{Тестовое покрытие всего проекта}
\label{tab:solver-coverage}
\end{table}

\end{document}
